% 00-cover.tex — Professional Cover Page + Preface (Anthropic-Style)
\thispagestyle{empty}
\begin{titlepage}
  \centering

  \vspace*{1.5cm}

  % --- Logo (TikZ decorative) ---
  \begin{tikzpicture}
    \node[
      circle, draw=accentClay, line width=1.5pt,
      fill=accentClay!8,
      minimum size=2.5cm,
      font=\sffamily\fontsize{28}{34}\selectfont\bfseries,
      text=accentClay
    ] {D};
  \end{tikzpicture}

  \vspace{1.2cm}

  % --- OpceanAI Brand Name (Large) ---
  {\sffamily\fontsize{44}{52}\selectfont\bfseries\color{inkDark} OpceanAI\par}

  \vspace{0.4cm}

  % --- Slogan in English ---
  {\sffamily\fontsize{13}{17}\selectfont\itshape\color{textgray}
    A sea of possibilities\par}

  \vspace{2.2cm}

  % --- Terracotta accent rule ---
  {\color{accentClay}\rule{0.5\textwidth}{1.2pt}}

  \vspace{1.5cm}

  % --- Report Title ---
  {\sffamily\fontsize{32}{38}\selectfont\bfseries\color{inkDark}
    Doki\par}

  \vspace{0.6cm}

  % --- Subtitle ---
  {\sffamily\fontsize{13}{17}\selectfont\color{textdark}
    A Complete Container Ecosystem\\for Android\par}

  \vspace{0.8cm}

  % --- Version ---
  {\sffamily\fontsize{11}{14}\selectfont\color{textgray}
    Technical Report --- Version 0.9.3\par}

  \vspace{3cm}

  % --- Author ---
  {\sffamily\fontsize{12}{16}\selectfont\color{inkDark}
    OpceanAI Contributors\par}

  \vspace{0.4cm}

  % --- Date ---
  {\sffamily\fontsize{10}{13}\selectfont\color{textgray} \today\par}

  \vfill

  % --- Bottom accent rule ---
  {\color{accentClay}\rule{\textwidth}{0.4pt}}

  \vspace{0.4cm}

  % --- Footer ---
  {\sffamily\fontsize{8}{10}\selectfont\color{textgray}
    Doki v0.9.3 \textbar{}
    \url{https://github.com/OpceanAI/Doki} \textbar{}
    Apache License, Version 2.0
  \par}

  \vspace{0.8cm}
\end{titlepage}

% PREFACE
\chapter*{Preface}
\addcontentsline{toc}{chapter}{Preface}

This technical report provides a comprehensive description of the Doki
container ecosystem, covering its architecture, components, and design
decisions. It is intended for developers, system administrators, and
researchers interested in container technology on Android platforms.

Doki is developed by OpceanAI and released as open-source software
under the Apache License, Version 2.0. The project began as a response
to the absence of robust containerization tools for Android
environments, where traditional solutions such as Docker are
unavailable due to kernel limitations and the lack of root access.

This report covers version 0.9.3 of Doki, which includes DokiLink-Lite
peer-to-peer mesh networking, 244 CLI commands, 12 isolation levels,
and support for both ARM64 and ARMv7 architectures.

\begin{center}
  \color{textgray}
  \sffamily\small
  OpceanAI \textbar{} \today
\end{center}

\cleardoublepage
